\begin{abstract}
Evaluating Large Language Models (LLMs) on repository-level feature implementation is a critical frontier in software engineering. However, establishing a benchmark that faithfully mirrors realistic development scenarios remains a significant challenge. Existing feature-level benchmarks generally suffer from two primary limitations: unrealistic task inputs enriched with code hints and significant data leakage risks due to their static nature. To address these limitations, we propose a new benchmark -- \textbf{FeatBench}, which introduces the following advances: \ding{182} \textbf{Realistic Task Inputs.} Task inputs consist solely of natural language requirements, strictly devoid of code hints (e.g., function signatures). This format mirrors realistic software development by requiring agents to independently bridge the gap between abstract user intent and concrete code changes. \ding{183} \textbf{Evolving Data.} FeatBench employs a fully automated pipeline to construct new benchmark versions from the latest repositories, effectively mitigating data contamination. The initial release comprises 157 tasks sourced from 27 actively maintained repositories. We evaluate two state-of-the-art agent frameworks with four leading LLMs on FeatBench. The results reveal that FeatBench poses a significant challenge, with the highest resolved rate reaching only 29.94\%. Crucially, our analysis uncovers a prevalent behavioral pattern of aggressive implementation, which leads to ``scope creep'' and widespread regressions where agents break existing features by diverging from the user's explicit intent. We release FeatBench, our automated pipeline, and all experimental results to facilitate further community research. Our code is available at \url{https://github.com/TsinghuaISE/FeatBench}.
\end{abstract}